
% Section 4 package for Paper 2
% Full spanning statement for the obstruction space

\begin{definition}
For each $d\ge 0$, let
\[
V_d^{(7)}=\Span_{\mathbb Q}\{n^kM_a(n):0\le k\le d,\ 1\le a\le 7\}.
\]
Let $E_d\subseteq \operatorname{ev}(V_d^{(7)})$ denote the Eisenstein-type prime-evaluation subspace, and define the obstruction space
\[
O_d:=\operatorname{ev}(V_d^{(7)})/E_d.
\]
For $f\in \operatorname{ev}(V_d^{(7)})$, we write $[f]$ for its image in $O_d$.
\end{definition}

\begin{lemma}\label{lem:Ed-polynomial}
Every element of $E_d$ is represented on primes by a polynomial in $p$.
\end{lemma}

\begin{proof}
By construction, $E_d$ is generated by prime restrictions of Eisenstein-type expressions arising from the divisor-sum parts of the basis elements $n^kM_a(n)$. For each divisor sum $\sigma_m(n)$, one has
\[
\sigma_m(p)=1+p^m
\]
at a prime $p$. Therefore every Eisenstein-type basis element evaluates on primes to a polynomial in $p$, and the same is true for every $\mathbb Q$-linear combination of such elements. Hence every element of $E_d$ is represented on primes by a polynomial in $p$.
\end{proof}

\begin{lemma}\label{lem:Qtau-not-in-Ed}
Let $Q\in\mathbb Q[p]$ be nonzero. Then
\[
Q(p)\tau(p)\notin E_d.
\]
Equivalently, the class $[Q(p)\tau(p)]$ is nonzero in $O_d$.
\end{lemma}

\begin{proof}
Suppose for contradiction that $Q(p)\tau(p)\in E_d$ for some nonzero polynomial $Q\in\mathbb Q[p]$. By Lemma~\ref{lem:Ed-polynomial}, there exists a polynomial $g\in\mathbb Q[p]$ such that
\[
Q(p)\tau(p)=g(p)
\]
for all primes $p$.

Since $Q$ is nonzero, it has only finitely many roots. Hence there exists $P_0$ such that
\[
Q(p)\neq 0
\qquad\text{for all primes }p>P_0.
\]
Therefore, for all primes $p>P_0$,
\[
\tau(p)=\frac{g(p)}{Q(p)}.
\]
The right-hand side is a nonzero rational function in $p$. Any nonzero rational function over $\mathbb Q$ has constant sign for all sufficiently large real inputs away from its poles. Since $Q$ has only finitely many zeros, it follows that $g(p)/Q(p)$ has constant sign for all sufficiently large primes.

Thus $\tau(p)$ would have constant sign for all sufficiently large primes. This contradicts the sign oscillation of Ramanujan's tau-function on primes, which follows from the Sato--Tate equidistribution of the normalized coefficients
\[
\frac{\tau(p)}{2p^{11/2}}.
\]
Hence no such polynomial $Q$ exists, and therefore
\[
Q(p)\tau(p)\notin E_d.
\]
\end{proof}

\begin{corollary}\label{cor:poly-tau-independent}
Let $m\ge 0$. Then the classes
\[
[\tau],\ [p\tau],\ \dots,\ [p^m\tau]
\]
are $\mathbb Q$-linearly independent in $O_d$.
\end{corollary}

\begin{proof}
Suppose
\[
a_0[\tau]+a_1[p\tau]+\cdots+a_m[p^m\tau]=0
\]
in $O_d$. By definition of the quotient, this means
\[
(a_0+a_1p+\cdots+a_mp^m)\tau(p)\in E_d.
\]
If the polynomial
\[
Q(p):=a_0+a_1p+\cdots+a_mp^m
\]
were nonzero, this would contradict Lemma~\ref{lem:Qtau-not-in-Ed}. Hence $Q=0$, and therefore
\[
a_0=a_1=\cdots=a_m=0.
\]
Thus the listed classes are linearly independent in $O_d$.
\end{proof}

\begin{proposition}\label{prop:Od-full-structure}
Assume the prime-level decompositions
\[
M_6(p)=f_6(p)+c_\tau\tau(p),
\qquad c_\tau\neq 0,
\]
and
\[
M_7(p)=f_7(p)+\alpha\tau(p)+\beta p\tau(p),
\qquad \beta\neq 0,
\]
with $f_6,f_7\in\mathbb Q[p]$. Then for every $d\ge 0$,
\[
O_d=\Span_{\mathbb Q}\{[\tau],[p\tau],\dots,[p^{d+1}\tau]\}.
\]
In particular,
\[
\dim O_d=d+2.
\]
\end{proposition}

\begin{proof}
We first prove spanning. Let
\[
\mathcal B_d^{(7)}=\{n^kM_a(n):0\le k\le d,\ 1\le a\le 7\}
\]
be the basis underlying $V_d^{(7)}$, and let $p^kM_a(p)$ denote the prime restriction of a basis element.

If $1\le a\le 5$, then $M_a(p)$ is polynomial in $p$, so $p^kM_a(p)\in E_d$. Hence these basis elements contribute nothing to $O_d$.

For $a=6$, one has
\[
p^kM_6(p)=p^k f_6(p)+c_\tau p^k\tau(p).
\]
Since $p^k f_6(p)$ is polynomial in $p$, its class in $O_d$ vanishes. Therefore
\[
[p^kM_6(p)]=c_\tau [p^k\tau(p)].
\]

For $a=7$, one has
\[
p^kM_7(p)=p^k f_7(p)+\alpha p^k\tau(p)+\beta p^{k+1}\tau(p).
\]
Again the polynomial term lies in $E_d$, so
\[
[p^kM_7(p)]=\alpha [p^k\tau(p)]+\beta [p^{k+1}\tau(p)].
\]

Thus every basis element of $\operatorname{ev}(V_d^{(7)})$ has class in
\[
\Span_{\mathbb Q}\{[\tau],[p\tau],\dots,[p^{d+1}\tau]\}.
\]
Hence
\[
O_d\subseteq \Span_{\mathbb Q}\{[\tau],[p\tau],\dots,[p^{d+1}\tau]\}.
\]

Conversely, because $c_\tau\neq 0$, the classes
\[
[\tau],[p\tau],\dots,[p^d\tau]
\]
occur through the elements
\[
[M_6(p)],\ [pM_6(p)],\ \dots,\ [p^dM_6(p)].
\]
Because $\beta\neq 0$, the class
\[
[p^{d+1}\tau]
\]
occurs through
\[
[p^dM_7(p)].
\]
Therefore
\[
\Span_{\mathbb Q}\{[\tau],[p\tau],\dots,[p^{d+1}\tau]\}\subseteq O_d.
\]

Combining the two inclusions gives
\[
O_d=\Span_{\mathbb Q}\{[\tau],[p\tau],\dots,[p^{d+1}\tau]\}.
\]

Finally, Corollary~\ref{cor:poly-tau-independent} shows that these $d+2$ classes are linearly independent. Hence
\[
\dim O_d=d+2.
\]
\end{proof}
